\chapter{Evaluation}
\label{ch:eval}

How do you know if your language model is any good?  In this chapter we cover \textbf{perplexity} and \textbf{bits per byte} (a tokenizer-independent metric), implement validation-loss tracking, and briefly survey zero-shot benchmarks and scaling laws.


\section{Perplexity}

Perplexity is the exponential of the average cross-entropy loss:
\[
  \text{PPL} = \exp\!\left(\frac{1}{N}\sum_{i=1}^{N} -\log p(x_i \mid x_{<i})\right) = e^{\,\mathcal{L}}
\]

Intuitively, perplexity measures how ``surprised'' the model is by the text.  A perplexity of 10 means the model is as uncertain as if it were choosing uniformly among 10 equally likely tokens at each position.

\textbf{Lower is better.}  A random model over a 32K vocabulary has perplexity $\sim$32\,768.  GPT-2 124M typically achieves perplexity $\sim$25--30 on web text.

\begin{lstlisting}[language=Python]
def perplexity(loss: float) -> float:
    import math
    return math.exp(loss)
\end{lstlisting}


\section{Bits Per Byte (BPB)}

Perplexity depends on the tokenizer: a model with a 50K vocabulary and a model with a 32K vocabulary are not directly comparable by perplexity because they predict over different token sets.

\textbf{Bits per byte} normalizes by the number of raw bytes in the text, making it tokenizer-independent:
\[
  \text{BPB} = \frac{\sum_i \mathcal{L}(x_i)}{\text{total bytes}} \cdot \frac{1}{\ln 2}
\]

Each token contributes its cross-entropy loss, weighted by how many UTF-8 bytes it represents.  The $1/\ln 2$ converts from nats (natural log) to bits.

BPB values for English text typically range from 0.8 (strong model) to 1.5 (weak model).  Random: $\sim$8 bits/byte.

\filetag{microgpt/eval.py}
\begin{lstlisting}
@torch.no_grad()
def compute_bpb(
    model: MicroGPT,
    val_loader,
    tokenizer: BPETokenizer,
    num_batches: int = 50,
) -> float:
    """Compute bits per byte (BPB) on validation data.

    BPB is tokenizer-independent: it measures compression in terms of
    raw bytes rather than tokens.  This allows fair comparison between
    models with different tokenizers.

    BPB = mean_cross_entropy / ln(2) * (tokens / bytes)
    """
    model.eval()
    total_loss = 0.0
    total_tokens = 0
    total_bytes = 0

    for i, (inputs, targets) in enumerate(val_loader):
        if i >= num_batches:
            break

        logits, _ = model(inputs)
        # Per-token cross-entropy
        B, T, V = logits.shape
        loss_per_token = F.cross_entropy(
            logits.reshape(-1, V),
            targets.reshape(-1),
            reduction="none",
        ).view(B, T)

        # Count bytes per token
        for b in range(B):
            for t in range(T):
                token_id = targets[b, t].item()
                try:
                    token_bytes = tokenizer.token_to_bytes(token_id)
                    n_bytes = len(token_bytes)
                except Exception:
                    n_bytes = 1  # fallback for special tokens

                total_loss += loss_per_token[b, t].item()
                total_bytes += n_bytes
                total_tokens += 1

    # BPB = total_nats / total_bytes / ln(2)
    import math
    bpb = total_loss / max(total_bytes, 1) / math.log(2)
    return bpb
\end{lstlisting}


\section{Validation Loss}

During training, we compute validation loss every \code{EVAL\_INTERVAL} steps on a held-out shard (the last parquet file).  This is the most important metric for detecting:

\begin{itemize}
  \item \textbf{Overfitting}: val loss starts increasing while train loss continues decreasing.  Unlikely at this data scale but important to monitor.
  \item \textbf{Training instability}: sudden spikes in val loss indicate exploding gradients or bad learning rate.
  \item \textbf{Convergence}: val loss plateaus, indicating diminishing returns from further training.
\end{itemize}


\section{Zero-Shot Benchmarks}

The GPT-2 paper's key claim was that language models can perform tasks zero-shot.  Standard benchmarks include:

\begin{center}
\begin{tabular}{lp{8cm}}
\toprule
Benchmark & What it tests \\
\midrule
\textbf{HellaSwag} & Commonsense reasoning: pick the most plausible continuation of a scenario. \\
\textbf{LAMBADA} & Long-range prediction: predict the last word of a passage. \\
\textbf{ARC} (AI2 Reasoning Challenge) & Science questions from grade-school exams. \\
\textbf{WinoGrande} & Pronoun resolution requiring commonsense. \\
\bottomrule
\end{tabular}
\end{center}

For zero-shot evaluation, the model is given each candidate answer and the one with the lowest perplexity (or highest log-probability) is selected.  No fine-tuning is performed.

A well-trained GPT-2 124M typically achieves:
\begin{itemize}
  \item HellaSwag: $\sim$30--35\% accuracy (random baseline: 25\%)
  \item LAMBADA: $\sim$40--50\% accuracy
\end{itemize}


\section{Scaling Laws}

Kaplan et al.\ (2020) and Hoffmann et al.\ (2022, ``Chinchilla'') established that language model loss follows a power law in compute:
\[
  L(C) \approx L_0 + A \cdot C^{-\alpha}
\]

Key findings:
\begin{itemize}
  \item Loss decreases predictably with more compute (larger models, more data, longer training).
  \item \textbf{Compute-optimal training}: for a fixed compute budget, there is an optimal ratio of model size to data size.  Chinchilla's ratio: $\sim$20 tokens per parameter.
  \item For our 135M parameter model, compute-optimal training would use $\sim$2.7B tokens.  Training on 10B tokens is ``over-training'' by Chinchilla standards, but smaller models benefit from seeing more data.
\end{itemize}

\tip{Nanochat reports a CORE score (Composite Overall Reasoning Evaluation) that aggregates multiple benchmarks.  Their GPT-2 reproduction achieves a CORE score of $\sim$0.257, matching the original GPT-2 124M.  This provides a useful single-number target for verifying your training run.}


\section{Summary}

Evaluating a language model:
\begin{enumerate}
  \item \textbf{Validation loss}: primary training metric; detect overfitting and instability.
  \item \textbf{Perplexity}: $e^{\text{loss}}$; intuitive but tokenizer-dependent.
  \item \textbf{Bits per byte}: tokenizer-independent compression metric; enables cross-model comparison.
  \item \textbf{Zero-shot benchmarks}: HellaSwag, LAMBADA, ARC test real capabilities.
  \item \textbf{Scaling laws}: predict performance at larger scale; guide compute allocation.
\end{enumerate}

The evaluation script is in \code{microgpt/eval.py}, runnable as:

\shelltag
\begin{lstlisting}[language=bash]
uv run python -m microgpt.eval \
    --checkpoint microgpt_checkpoints/best.pt \
    --tokenizer tokenizer
\end{lstlisting}
